\documentclass[10pt,a4paper]{amsart}
\usepackage{amsmath,amssymb,amsthm,graphicx}
\RequirePackage{amssymb,amsmath,amsthm,bm}
\RequirePackage{thmtools}
\RequirePackage[capitalize]{cleveref}
\providecommand{\tightlist}{%
  \setlength{\itemsep}{0pt}\setlength{\parskip}{0pt}}

%%%%% varphi instead of phi
\renewcommand{\phi}{\varphi}

\mathchardef\mhyph="2D 
\providecommand{\A}{\mathcal{A}}
\providecommand{\C}{\mathcal{C}}
\providecommand{\L}{\mathcal{L}}
\providecommand{\B}{\mathcal{B}}
\providecommand{\M}{\mathcal{M}}
\renewcommand{\S}{\mathcal{S}}
\providecommand{\ol}[1]{\overline{#1}}
\providecommand{\mbf}[1]{\mathbf{#1}}
\providecommand{\mf}[1]{\mathfrak{#1}}
\providecommand{\mrm}[1]{\mathrm{#1}}
\providecommand{\mc}[1]{\mathcal{#1}}
\providecommand{\ra}{\rightarrow}
\providecommand{\Ra}{\Rightarrow}
\providecommand{\La}{\Leftarrow}
\providecommand{\lr}{\leftrightarrow}
\providecommand{\hf}{\mathbb{HF}}
\providecommand{\spec}{DgSp}
\providecommand{\bSigma}{\bm\Sigma}
\providecommand{\bDelta}{\bm\Delta}
\providecommand{\bPi}{\bm\Pi}
\providecommand{\Pinf}[1]{\Pi^{\mathrm{in}}_{#1}}
\providecommand{\Sinf}[1]{\Sigma^{\mathrm{in}}_{#1}}
\providecommand{\dSinf}[1]{d\mhyph\Sigma^{\mathrm{in}}_{#1}}
\providecommand{\Dinf}[1]{\Delta^{\mathrm{in}}_{#1}}
\providecommand{\bp}{\bar p}
\providecommand{\bq}{\bar q}
\providecommand{\ba}{\bar a}
\providecommand{\bb}{\bar b}
\providecommand{\bc}{\bar c}
\providecommand{\bd}{\bar d}
\providecommand{\bx}{\bar x}
\providecommand{\by}{\bar y}
\providecommand{\bz}{\bar z}
\providecommand{\LO}{\mathrm{LO}}
\declaretheorem[name={Theorem},numberwithin=section]{theorem}
\declaretheorem[name={Proposition}, sibling=theorem]{proposition}
\declaretheorem[name={Lemma}, sibling=theorem]{lemma}
\declaretheorem[name={Fact}, sibling=theorem]{fact}
\declaretheorem[name={Claim}, numberwithin=theorem]{claim}
\declaretheorem[name={Definition},numberwithin=section,sibling=theorem, style=definition]{definition}
\declaretheorem[name={Remark}, numbered=no,style=definition]{remark}
\declaretheorem[name={Corollary}, sibling=theorem]{corollary}
\declaretheorem[name={Question},sibling=theorem]{question}
\newcommand{\T}{{\mathcal{T}}}
\newcommand{\SR}{SR}
\newcommand{\pSR}{pSR}
\newcommand{\Q}{\mathbb{Q}}
\newcommand{\Z}{\mathbb{Z}}
\newcommand{\K}{\mathcal{K}}
\newcommand{\N}{\mathcal{N}}
\renewcommand{\L}{\mathcal{L}}
\providecommand{\Ginf}[1]{\Gamma^{\mathrm{in}}_{#1}}
\title{Comparing tuple sets in the labelled tree reduces to comparing the tuples coordinatewise}
\author{David Gonzalez, Dino Rossegger}
\date{arxiv:2305.07126v1; CC BY 4.0}
\begin{document}
\maketitle
\noindent\textit{Editorial note.} The counted statement is the author environment at source0.tex lines 509-513 and its proof environment opens at line 516, the line immediately after; the pairing is by adjacency and the retained proof range 516-544 is the whole environment.\par\medskip
\noindent\textit{Editorial note.} The retained preamble is the author style file dino.sty (package lines, the theorem declaration support, the notation and tuple macros and the theorem environments) followed by the ordinal and rank macros of the manuscript itself, reproduced verbatim; this is the same recipe as the previously published card of this article, so no package outside it is requested and extra\_packages stays empty.\par\medskip
\noindent\textit{Editorial note.} Metadata provenance: source\_id and source\_version are the values recorded in this source licence evidence file (sources/scott-linear-orders-2305.07126v1/LICENSE\_EVIDENCE.json: arxiv:2305.07126v1 / 2305.07126v1), which the card points at directly; the title and author list are transcribed from the authors\_index and title keys of the source record of the pinned copy. The card carries no counter setting and imposes no original numbering, since this article has no label-number dump and its numbering cannot be reconstructed from the retained ranges alone.\par\medskip
\noindent\textit{Editorial note.} No context span and no external reference are needed: the retained statement and proof use no label or citation from outside their own ranges.\par\medskip
\noindent\textit{Editorial scope.} One original result of the pinned author manuscript, counted once: the coordinatewise criterion for the back-and-forth comparison of isomorphic downward closed subsets of the labelled tree (author label lemma:factoring).\par\medskip
\noindent\textit{Allowed standard prerequisites (editorial).}\begin{itemize}
\item Back-and-forth relations on the labelled tree. The relations <=\_alpha on tuples of the labelled tree are the usual back-and-forth relations of rank alpha, and they are monotone.. They are the relation being characterised.
\item Isomorphic downward closed subsets of the tree. S and T are finite downward closed subsets of the labelled tree, enumerated in an isomorphism.. They are the two sides compared.
\item Infinitely replicated branches. Every node of the labelled tree has infinitely many copies of its branch above it.. It makes the upper cones of a node and of its class isomorphic.
\end{itemize}\medskip
\begin{lemma}\label{lemma:factoring}
Let $S=(x_1,\cdots,x_n)$ and $T=(y_1,\cdots,y_n)$ be (finite) isomorphic
downward closed subsets of $\T_\A$.
$S\leq_\alpha T$ if and only if for all $i$, $x_i\leq_\alpha y_i$. 
\end{lemma}

\begin{proof}
If $S\leq_\alpha T$, $x_i\leq_\alpha y_i$ for any $i$ follows immediately from the monotonicity of $\leq_\alpha$.


Say that for all $i$, $x_i\leq_\alpha y_i$.
For $z\in \T_\A$ Let $U(z)$ be all elements above $z$ and note that
for all $i$, $U(x_i)\leq_\alpha U(y_i)$.
This holds because any $\Pinf{\alpha}$ formula about $U(x_i)$ can be transformed into a $\Pinf{\alpha}$ formula about $x_i$ by restricting all quantifiers to being above $x_i$.

Let $U_S(x_i) = \{y\in T_A\vert y\geq x_i \land \lnot\exists j ~ x_j>x_i \land y\geq x_j\}$ and define $U_T(y_i)$ analogously.
Because $S$ is downward closed, the $U_S(x_i)$ partition $\T_\A$ into disjoint sets (and similarly for $T$).
Because the branches of $T_A$ are infinitely replicated, we can observe that $U(x_i)\cong U_S(x_i)$ and $U(y_i)\cong U_T(y_i)$.
In particular, this means that for all $i$, $U_S(x_i)\leq_\alpha U_T(y_i)$.

For the sake of contradiction, say that $S\not\leq_\alpha T$.
(Note that $\alpha\neq0$, as $S$ and $T$ are isomorphic.)
This means there is a $\beta<\alpha$ and $\bd\in T_A$ such that for all $\bc\in T_A$
$(S,\bc)\not\geq_\alpha (T,\bd)$.
Write $\bd=(\bd_1\cdots\bd_n)$ when $\bd_i$ is the (possibly empty) subset of
$\bd$ that is in $U_T(y_i)$. For every $i$, let $\bc_i$ be the $\exists$-player's winning response to $\bd_i$ in $U_S(x_i)$ and let $\bc=(\bc_1\cdots \bc_n)$.
This is a valid move because the elements in  $U_T(y_i)$ and $U_T(y_j)$ for distinct $i$ and $j$ are never comparable to each other. 
As the game proceeds, the $\exists$-player may continue to play according to this strategy for the same reason.
When the game reduces to the $0$ level, it is assured that the $\exists$-player has won on all of the $U_S(x_i)$ and  $U_T(y_i)$.
In other words, the relations of the final tuple within these structures is the same.
As observed above, there are never any relations between points in distinct $U_S(x_i)$ or  $U_T(y_i)$.
Therefore, the relations between the structures are also the same.
Thus we have described a winning strategy for the $\exists$-player contradicting
that $S\not\leq_\alpha T$
\end{proof}

\end{document}
